\chapter{Incident Response Workflow}
\label{ch:incident-response-workflow}

This chapter covers the incident response lifecycle as implemented by the \texttt{ops.incident}
wrapper. It walks through declaring an incident, commander takeover, triage, escalation,
and resolution, concluding with a complete annotated example.

% ─────────────────────────────────────────────────────────────────────────────
\section{Incident Lifecycle}
\label{sec:ir-lifecycle}
% ─────────────────────────────────────────────────────────────────────────────

\texttt{gert.ops} models the incident lifecycle as six stages:

\begin{description}
  \item[Detect] An alert fires or a human observes symptoms. The on-call engineer is paged.
  \item[Declare] The on-call engineer starts the incident runbook. The \texttt{ops.incident}
    wrapper activates, declaring the incident and assigning an Incident Commander.
  \item[Triage] Responders gather data, run diagnostics, and determine root cause or impact
    scope. Each triage step requires evidence submission.
  \item[Mitigate] Immediate steps to reduce customer impact (e.g., restart a service, roll
    back a deploy, apply a hotfix). Mitigation is tracked against the SLA timeout.
  \item[Resolve] The root cause is addressed and the service is fully restored.
    The Incident Commander signs off on resolution.
  \item[Review] Post-incident review (PIR). The evidence bundle from the run is the
    primary artifact for the PIR meeting. This stage is outside the runbook scope.
\end{description}

% ─────────────────────────────────────────────────────────────────────────────
\section{Declaring an Incident}
\label{sec:ir-declare}
% ─────────────────────────────────────────────────────────────────────────────

Starting a runbook that contains an \texttt{ops.incident} wrapper immediately writes an
\texttt{incident/declared} trace event and transitions the run into incident mode. In
incident mode:

\begin{itemize}
  \item All \texttt{ops.manual} steps with \texttt{requires\_role: responder} are unlocked
    for any user listed in \texttt{meta.roles.responder}.
  \item SLA timers start immediately. The \texttt{mitigation-timeout} timer tracks time to
    first mitigation. The \texttt{resolution-timeout} timer tracks total incident duration.
  \item The declared \texttt{commander} assumes DRI authority for all steps within the
    \texttt{ops.incident} wrapper, even if a different DRI is declared at the runbook level.
\end{itemize}

% ─────────────────────────────────────────────────────────────────────────────
\section{Incident Commander Takeover}
\label{sec:ir-commander-takeover}
% ─────────────────────────────────────────────────────────────────────────────

The Incident Commander (IC) may take over at any point during the run. Takeover is
supported in two ways:

\paragraph{Declared takeover.} The IC is declared in \texttt{ops.incident.commander}. When
the wrapper starts, the IC immediately holds DRI authority.

\paragraph{Runtime takeover.} During a live run, the IC can execute:

\begin{verbatim}
$ gert incident take-command <run-id> --as=@ic-oncall
\end{verbatim}

This writes an \texttt{incident/commanderTakeover} trace event and reassigns DRI authority.
The previous DRI is notified and transitions to \texttt{observer} status for subsequent steps.

\paragraph{DRI transfer back.} After the incident is resolved, the IC can transfer DRI back
to the original owner:

\begin{verbatim}
$ gert incident transfer-command <run-id> --to=@alice
\end{verbatim}

% ─────────────────────────────────────────────────────────────────────────────
\section{Triage Steps}
\label{sec:ir-triage}
% ─────────────────────────────────────────────────────────────────────────────

Triage steps use \texttt{ops.manual} with \texttt{requires\_role: responder}. Each triage step
should:
\begin{itemize}
  \item Provide specific, actionable instructions (dashboards, runbook links, check commands).
  \item Require an attestation evidence record summarizing the findings.
  \item Use \texttt{ops.cli} sub-steps (or direct \texttt{ops.cli} steps) for diagnostic
    commands, capturing output as evidence.
\end{itemize}

\paragraph{Branch on triage outcome.} Use the core \texttt{branch} step type within the
incident wrapper to route to different mitigation steps based on triage findings.

% ─────────────────────────────────────────────────────────────────────────────
\section{Escalation When SLA Timeout Fires}
\label{sec:ir-escalation}
% ─────────────────────────────────────────────────────────────────────────────

When the \texttt{mitigation-timeout} fires, gert:
\begin{enumerate}
  \item Writes an \texttt{incident/slaBreached} trace event with the timeout type and elapsed
    duration.
  \item Notifies the \texttt{escalation-target} with the run ID, current step, and elapsed time.
  \item The run continues. The escalation is informational; it does not pause or abort the run.
\end{enumerate}

When the \texttt{resolution-timeout} fires, gert:
\begin{enumerate}
  \item Writes an \texttt{incident/slaBreached} trace event.
  \item Notifies the \texttt{escalation-target}.
  \item Suspends the run and requires IC acknowledgement to continue.
\end{enumerate}

% ─────────────────────────────────────────────────────────────────────────────
\section{Complete Example: Service Degradation Incident}
\label{sec:ir-complete-example}
% ─────────────────────────────────────────────────────────────────────────────

The following runbook covers a service degradation incident from initial page to resolution.

\begin{minted}{yaml}
apiVersion: runbook/v2
kit: ops/v1

meta:
  name: service-degradation-response
  kind: incident
  description: >
    Incident response playbook for service degradation alerts.
    Covers triage, mitigation (rollback or restart), and resolution.
  inputs:
    service_name:
      from: prompt
      description: "Affected service (e.g., payment-processor)"
    incident_ticket:
      from: prompt
      description: "Incident ticket ID (e.g., INC-20240514-001)"
    alert_description:
      from: prompt
      description: "Brief description of the alert that triggered this runbook"
  roles:
    dri: "@oncall-sre"
    responder:
      - "@oncall-sre"
      - "@oncall-backend"
    incident-commander: "@ic-oncall"
    observer:
      - "@sre-team"
      - "@management-alerts"
  sla:
    mitigation-timeout: 30m
    resolution-timeout: 4h
    escalation-target: "@vp-engineering"
  governance:
    allowed_commands: [kubectl, curl, jq, pg_isready, redis-cli, stern]
    deny_env_vars: ["*_SECRET*", "*_TOKEN", "*_KEY"]

flow:
  # ── Incident declaration and lifecycle wrapper ─────────────────────────────
  - step:
      id: incident_response
      type: ops.incident
      title: "Service degradation: {{ .service_name }}"
      severity: sev2
      incident-id: "{{ .incident_ticket }}"
      commander: "@ic-oncall"
      sla:
        mitigation-timeout: 30m
        resolution-timeout: 4h
        escalation-target: "@vp-engineering"
      steps:
        # ── Stage 1: Initial triage ──────────────────────────────────────────
        - step:
            id: initial_triage
            type: ops.manual
            title: "Initial triage"
            requires_role: responder
            instructions: |
              Alert: {{ .alert_description }}
              Service: {{ .service_name }}

              Check the following:
              1. Error rate: https://grafana.internal/d/{{ .service_name }}/errors
              2. Latency: https://grafana.internal/d/{{ .service_name }}/latency
              3. Recent deployments in the last 2 hours:
                 kubectl rollout history deployment/{{ .service_name }} -n production
              4. Pod health:
                 kubectl get pods -n production -l app={{ .service_name }}
              5. Recent logs (last 100 lines):
                 stern {{ .service_name }} -n production --tail=100
            evidence:
              - type: ops.evidence.attestation
                label: "Initial triage findings"
                required: true
                prompt: >
                  Describe the symptoms: error rate, latency, pod status, and
                  whether a recent deployment is likely the cause.

        # ── Stage 2: Diagnostic commands ─────────────────────────────────────
        - step:
            id: check_pod_status
            type: ops.cli
            title: "Check pod status"
            command: kubectl
            args:
              - get
              - pods
              - "-n"
              - production
              - "-l"
              - "app={{ .service_name }}"
              - "-o"
              - wide
            capture:
              pod_status: stdout
            evidence:
              auto: true
              label: "Pod status"

        - step:
            id: check_recent_events
            type: ops.cli
            title: "Check recent Kubernetes events"
            command: kubectl
            args:
              - get
              - events
              - "-n"
              - production
              - "--sort-by=.lastTimestamp"
              - "--field-selector"
              - "involvedObject.name={{ .service_name }}"
            capture:
              k8s_events: stdout
            evidence:
              auto: true
              label: "Recent Kubernetes events"

        # ── Stage 3: Mitigation decision ─────────────────────────────────────
        - step:
            id: mitigation_decision
            type: ops.manual
            title: "Select mitigation strategy"
            requires_role: incident-commander
            instructions: |
              Pod status: {{ .pod_status }}
              Events: {{ .k8s_events }}

              Based on the triage findings, choose the mitigation strategy:
              - If a recent deployment is the cause: select ROLLBACK
              - If pods are crash-looping with no recent deploy: select RESTART
              - If the issue is external (DB, upstream): select ESCALATE

              Document your reasoning before selecting.
            evidence:
              - type: ops.evidence.attestation
                label: "IC mitigation decision"
                required: true
                prompt: >
                  State the root cause hypothesis and chosen mitigation strategy.
                  Justify why this strategy was selected.

        # ── Stage 4a: Rollback mitigation ────────────────────────────────────
        # (This branch is executed when the IC determines rollback is needed.
        #  In practice, wrap in a branch step on the mitigation_decision choice.)
        - step:
            id: rollback_deployment
            type: ops.cli
            title: "Rollback to previous deployment"
            command: kubectl
            args:
              - rollout
              - undo
              - "deployment/{{ .service_name }}"
              - "-n"
              - production
            capture:
              rollback_output: stdout
            evidence:
              auto: true
              label: "kubectl rollout undo output"

        - step:
            id: wait_rollback
            type: ops.cli
            title: "Wait for rollback to complete"
            command: kubectl
            args:
              - rollout
              - status
              - "deployment/{{ .service_name }}"
              - "-n"
              - production
              - "--timeout=5m"
            evidence:
              auto: true
              label: "rollback rollout status"

        # ── Stage 5: Post-mitigation health check ────────────────────────────
        - step:
            id: mitigation_health_check
            type: ops.cli
            title: "Check service health endpoint"
            command: curl
            args:
              - "-sf"
              - "https://{{ .service_name }}.prod.internal/healthz"
              - "-o"
              - /dev/null
              - "-w"
              - "%{http_code}"
            capture:
              health_code: stdout
            evidence:
              auto: true
              label: "health endpoint response code"

        # ── Stage 6: Resolution sign-off ──────────────────────────────────────
        - step:
            id: resolution_sign_off
            type: ops.manual
            title: "IC resolution sign-off"
            requires_role: incident-commander
            instructions: |
              Confirm all of the following before signing off:
              1. Error rate has returned to normal (< 0.1%).
              2. Latency is within SLO (p99 < 200ms).
              3. No crash-looping pods.
              4. Health endpoint returns HTTP 200.

              Dashboard: https://grafana.internal/d/{{ .service_name }}
            evidence:
              - type: ops.evidence.screenshot
                label: "Post-mitigation Grafana dashboard"
                required: true
                description: "Dashboard showing error rate and latency, 10 min post-mitigation."
              - type: ops.evidence.attestation
                label: "IC resolution declaration"
                required: true
                prompt: >
                  I declare incident {{ .incident_ticket }} resolved.
                  {{ .service_name }} is healthy. Mitigation was successful.
                  Root cause: [state root cause].
                  Action items for PIR: [list items or 'none'].
\end{minted}
